% Figure source: addition and concatenation of token and positional vectors.
% Uses the standard LaTeX picture environment; no additional package is needed.

\begin{figure}[htbp]
    \centering
    \setlength{\unitlength}{1mm}
    \resizebox{\textwidth}{!}{%
    \begin{picture}(174,64)
        \put(0,59){\makebox(0,0)[l]{\textbf{Addition}}}
        \put(0,42){\framebox(29,10){\shortstack{Token embedding\\$\mathbf{x}_i\in\mathbb{R}^{d}$}}}
        \put(0,15){\framebox(29,10){\shortstack{Position vector\\$\mathbf{p}_i\in\mathbb{R}^{d}$}}}
        \put(29,47){\line(1,-1){11}}
        \put(29,20){\line(1,1){11}}
        \put(43,33){\makebox(0,0){\Large \(+\)}}
        \put(47,33){\vector(1,0){11}}
        \put(58,27){\framebox(35,12){\shortstack{Combined vector\\$\mathbf{x}_i+\mathbf{p}_i$}}}

        \put(104,59){\makebox(0,0)[l]{\textbf{Concatenation}}}
        \put(104,42){\framebox(27,10){\shortstack{Content\\$\mathbf{x}'_i$}}}
        \put(104,15){\framebox(27,10){\shortstack{Position\\$\mathbf{p}'_i$}}}
        \put(131,47){\line(1,-1){10}}
        \put(131,20){\line(1,1){10}}
        \put(144,33){\makebox(0,0){\Large \(\oplus\)}}
        \put(148,33){\vector(1,0){8}}
        \put(156,27){\framebox(18,12){\shortstack{Joined\\vector}}}
    \end{picture}%
    }
    \caption{Two ways of incorporating positional information into token
    representations. Addition superimposes full-dimensional vectors,
    whereas concatenation joins half-dimensional content and positional
    components while retaining the same total dimension. Adapted
    conceptually from the comparison in \textcite{zhou2024rethinking}.}
    \label{fig:pe-addition-concatenation}
\end{figure}
